%% acknowledgments.tex

% From mitthesis package
% Version: 1.02, 2024/06/19
% Documentation: https://ctan.org/pkg/mitthesis
% Author: Alejandro Santomé Valverde

\chapter*{Acknowledgments}
\pdfbookmark[0]{Acknowledgments}{acknowledgments}


Embarking on an industrial PhD is akin to walking a tightrope between the profound, meticulous 
rigor of academia and the relentless, pragmatic pace of industry. It is a journey that cannot be 
navigated alone, and looking back at the winding path that led to this manuscript, I realize how 
fortunate I have been in my traveling companions.

My deepest gratitude belongs to my supervisors, Dr. Daniel Pérez and Prof. José Capmany, who 
provided two vital and complementary pillars for this journey. A special word of appreciation goes 
to Dani. You first introduced me to the iPronics vision with a contagious passion that convinced me 
to embark on this adventure. Witnessing your evolution into the CTO of a massively growing company 
has been extraordinary, and I am deeply thankful not only for your leadership but also for the time 
and effort dedicated to meticulously reviewing this manuscript. I also want to express my sincere 
gratitude to Prof. Capmany; while he was less involved in the daily industrial grind, his stature 
and legacy ensured I never lost sight of the fundamental academic rigor this thesis demands. I am 
equally grateful to my thesis reviewers, Prof. Mahdi Nikdast and Prof. Nikos Pleros, for their 
insightful feedback and the valuable time dedicated to evaluating this work.

This journey is inextricably linked to iPronics, the organization that provided the framework, 
resources, and trust to make this industrial PhD possible. I joined as one of its first ten 
employees, and none of this chapter would exist without my colleague and friend, Mikel Gutiérrez. 
You were there when the team was just four people and opened the door to this adventure; watching 
your immense professional growth ever since has been a constant source of inspiration to me. In 
those early startup days, survival and growth meant wearing every possible hat, from designing 
photonic chips in the morning and stripping wires to assemble equipment in the afternoon, granting 
me an invaluable transversal vision. During those formative years, I was fortunate to learn from 
our validation manager, Jesús Benítez, who patiently taught me the intricate arts of the 
laboratory—an education that Ana Gutiérrez subsequently helped to refine.

As the journey progressed, my path pivoted towards the software department—a transition that 
occurred at a time when I severely doubted my capacity to finish this thesis. It was then that Paco 
Brevers, the software manager, stepped in to help me navigate that uncertainty. He did not 
intervene with academic advice or technical content, but rather with something far more crucial: 
unconditional support and the exact space I needed to execute my work. He gave me the resources, 
the time, and the leverage to ask for whatever was necessary to cross the finish line—not because 
it was critical for the department's roadmap, but simply because he knew it was deeply important to 
me. That selfless backing was indispensable, and I will be profoundly grateful to him for the rest 
of my life. Within this team, I also experienced immense growth in true collaboration. I had the 
privilege of learning from top-tier technical professionals who possess the rare quality of 
dropping whatever they are doing to help a colleague in need, truly embodying a 'team first, 
individual second' mindset. Alejandro Cruz stands as an absolute pillar of this philosophy. He is 
that indispensable person everyone turns to when they hit a wall; perhaps with a bit of customary 
grumbling, he always steps up and finds the solution. iPronics is extraordinarily lucky to have him 
as its employee number one.

A doctoral journey is a unique endeavor best understood by those walking the same arduous path. I 
want to share this milestone with my fellow PhD candidates: David, even though our daily paths at 
iPronics have diverged, Vicente (Zhenyun Xie), and Cristina. We shared the struggles, the doubts, 
and the triumphs. I must also extend my sincere gratitude to Luis, the manager of the Photonics 
department, who acted as a crucial mentor in shaping and refining the final stages of this 
manuscript. To the entire Photonics team under his wing, and to the relentless Validation and Lab 
team: thank you for laying the physical groundwork that made the science possible. 

Within our software ecosystem, I want to highlight the Calibration and 'Core Software' teams, 
masterfully led by Mikel. Erica, your support was a key catalyst in finalizing this thesis, and 
working with you is a pleasure. I cannot forget the 'software dark side', who taught us to elevate 
our collaborative standards. On a transversal note, I must thank Carlos López from the Projects 
team. His arrival provided a true captain for the Prometheus and Allegro European projects, 
bringing a focus that directly aided our results.


Every journey, no matter how far it takes us professionally, is sustained by our roots. Therefore, 
I would like to switch to my mother tongue to thank those who provided the very foundation of this 
path.

Todo camino tiene un punto de partida, y el mío siempre ha estado cimentado en mi familia. A mis 
padres, gracias por darme la oportunidad incalculable de estudiar, de salir ahí afuera, de ver 
mundo y de nutrirme de conocimiento. Vengo de una familia sin tradición de estudios superiores, y, 
aun así, confiasteis ciegamente en que mi vocación era la correcta. Me apoyasteis sin dudarlo, 
incluso cuando el esfuerzo económico no era nada sencillo. Con vuestros caracteres tan distintos y 
vuestras formas de demostrar afecto que no se parecen en nada, siempre habéis estado ahí, 
sosteniendo cada uno de mis pasos.

Este viaje también ha estado acompañado por quienes me dan aire fuera del laboratorio: mis amigos 
de toda la vida de Moaña, mis compañeros de la Universidad de Vigo y mis amigos de Valencia, que 
desde el primer día me acogieron como a uno más. Dentro de este grupo, necesito hacer una mención 
muy especial a mi mejor amiga, Laura. Aunque siempre deseas que vuelva a casa, nunca dejaste de 
animarme a terminar este doctorado. Entendiste lo importante que era para mí y, a pesar de la 
distancia, sé que te alegras de corazón al verme feliz aquí.

Finalmente, el cierre de este camino te pertenece a ti, Lau. Empezaste esta historia a mi lado 
siendo mi novia y la terminamos siendo marido y mujer. Contigo comparto los logros y las penas, lo 
bueno y lo malo. Eres quien hace que los días de celebración sean mucho mejores y quien consigue 
suavizar los días más oscuros. Eres una verdadera inspiración para mí, tanto a nivel personal como 
profesional. He visto cómo la vida te ha puesto muchas trabas y cómo siempre logras salir adelante 
sin perder un solo ápice de tu esencia. Eres la mejor persona que he conocido; tengo una suerte 
inmensa de que me hayas acompañado durante toda esta tesis y, por encima de todo, de que seas mi 
mujer.
